\documentclass{article}

\usepackage[english]{babel}

\usepackage[letterpaper,top=2cm,bottom=2cm,left=3cm,right=3cm,marginparwidth=1.75cm]{geometry}

\usepackage{amsmath}
\usepackage{graphicx}
\usepackage[strings]{underscore}
\usepackage[colorlinks=true, allcolors=blue]{hyperref}
\usepackage{float}

\setlength{\parindent}{0pt}

\title{Distributed Systems Exercise 2 Report}

\author{
    Aditya Prashantrao Kapse \and
    Lukas Linss \and
    Michal \and
    Navina \and
    Pranjal \and
    Shravani Chougule
}

\begin{document}
\maketitle
\tableofcontents
\newpage

\section{Requirements}
\subsection{Problem Context and Scope}

The system being designed is a globally accessible journey pre-booking service. Every road vehicle driver in the world must use this service to request authorisation for any journey they wish to make. No driver may begin a journey without having received a notification confirming that their requested journey is acceptable. The system must therefore be highly available, correct under concurrent load, and capable of scaling to serve millions of users distributed across every timezone and geographic region. During peak commute hours, these requests will arrive in bursts. The system must handle these bursts without degrading in correctness or availability.\\
At the same time, the instructions explicitly warn against over-engineering. A system that provides perfect consistency everywhere, with synchronous cross-continent coordination on every booking, would be too slow and too expensive to operate at global scale. The design must therefore make deliberate trade-offs, being strict where correctness is critical and relaxed where it is safe to do so.

\subsection{Assumptions About the Failure Model}
Before specifying requirements, it is important to be explicit about what kinds of failures the system is designed to tolerate. This is called the failure model, and every design decision that follows is made with respect to it.\\

\textbf{Crash-recovery failures:} The system assumes that services and database nodes fail by crashing and then eventually recovering, rather than by behaving maliciously or sending incorrect data (Byzantine failures). This is the standard assumption for distributed systems built on commodity cloud infrastructure. A crashed service may lose any in-memory state it had not yet persisted, but it will not send corrupt or deliberately wrong messages.
\\

\textbf{Network partitions:} The system assumes that network partitions can and do occur, particularly between geographically separated data centres. During a partition, some nodes may be unable to communicate with others. The system must continue operating (possibly in a degraded mode) on both sides of a partition and converge to a consistent state when the partition heals.
\\

\textbf{Partial failures:} Individual services or database replicas may fail independently. The system should not treat the failure of one component as a reason to bring down the entire system. Services that are not directly involved in a failing operation should continue to function normally.\\

\textbf{Failure frequency:} Hardware failures in a cloud environment are rare on a per-node basis but become statistically frequent at the scale of hundreds or thousands of nodes. Network latency spikes and brief connectivity losses between availability zones are common. Full data centre outages are rare but must be planned for. The system is therefore designed to tolerate individual node failures routinely, availability zone failures occasionally, and full region failures as a worst-case scenario.\\
\textbf{Data loss:} The system assumes that committed transactions in the database are durable — that is, once a booking has been acknowledged to a driver as received, it will not be silently lost. This requires write-ahead logging and persistent storage on all database nodes.

\subsection{Functional Requirements}
The functional requirements describe what the system must do from the perspective of its users and operators.\\

\textbf{Driver-facing requirements:}\\
The system must allow drivers to register an account with their personal details, driving licence number, vehicle type, and home region. Once registered, a driver must be able to authenticate and receive a session token that they can use for all subsequent interactions.\\
A driver must be able to look up a route between any two geographic coordinates anywhere in the world. The system must return the sequence of road segments that make up the route, along with an indication of current segment utilisation so the driver can make an informed choice about departure time.\\
A driver must be able to submit a booking for a specific route and departure time. The booking request must be acknowledged immediately, the driver should not be made to wait for the full capacity check to complete before receiving confirmation that their request has been received. The driver must be notified of the outcome (approved or rejected) through an inbox notification as soon as the system has made its determination.\\
A driver must be able to view their current and past bookings, cancel a pending booking that has not yet been resolved, and view their inbox messages.\\

\textbf{Enforcement support requirements:}\\
The system must ensure that no booking is approved if doing so would cause a road segment to exceed its predefined capacity during the booking's time window. This check must be atomic under concurrent load — two drivers competing for the last available slot on a segment must not both be approved.\\
The system must expose booking status information in a form that external enforcement systems (such as roadside cameras or checkpoint systems) can query. A booking must be unambiguously in one of a small number of states: PENDING, APPROVED, REJECTED, or EXPIRED. An enforcement system can query whether a given driver has an APPROVED booking for their current location and time without needing to understand the internal processing pipeline.\\
The system must automatically expire bookings whose departure time passes without resolution, so that stale reservations do not permanently consume segment capacity.

\subsection{Non-Functional Requirements}
Non-functional requirements describe the performance, reliability, and operational characteristics the system must meet. These are not arbitrary, each one is derived from the expected usage and failure patterns of the system.\\
\textbf{Performance:}\\ Booking submission must return an acknowledgement within 200 milliseconds for drivers. This requires regional deployment so that no driver is more than one or two network hops from a service endpoint. Route lookups must return within 500 milliseconds for cached routes and within 2 seconds for routes that require fresh computation. Notification delivery must complete within 5 seconds of a booking being assessed.\\

\textbf{Scalability:}\\
The system must be designed from the outset to scale horizontally. No single component may be a bottleneck that cannot be relieved by adding more instances. At steady state, the system must be capable of handling at least 100,000 booking requests per second globally. Individual services must be independently scalable so that, for example, the conflict detection service can be scaled up during peak hours without unnecessarily scaling the messaging service. \\

\textbf{Availability:}\\ The booking submission and status query endpoints must achieve 99.9\% availability (no more than approximately 8.7 hours of downtime per year). The system must remain available in the event of a single availability zone failure within a region. Planned maintenance must be possible without taking the system offline. \\

\textbf{Reliability:} No acknowledged booking request must ever be silently lost. If the system accepts a booking submission and returns an acknowledgement to the driver, that booking must be processed to a terminal state (APPROVED, REJECTED, or EXPIRED) even if multiple services subsequently fail and recover.\\

\textbf{Consistency:} Segment capacity must never be overbooked. This is a hard correctness requirement that cannot be relaxed. All other consistency requirements (notification delivery, booking status visibility) may be eventually consistent — there is an acceptable window of a few seconds during which a booking status may have changed but the driver has not yet been notified.\\

\textbf{Durability:} All committed booking data must survive the simultaneous loss of any single node. In production deployments, data must survive the loss of an entire availability zone. \\

\textbf{Security:} All client-facing API endpoints must require authentication. Drivers must only be able to read and modify their own bookings. Rate limiting must be enforced to prevent individual drivers from overwhelming the system.

\subsection{Usage Pattern Analysis}
Understanding the expected usage patterns is essential for making good design decisions. The following observations shape the architecture significantly.\\

\textbf{Read-heavy workload:}\\ The vast majority of interactions with the system are reads. Drivers check their booking status, browse routes, and read their inbox far more frequently than they submit new bookings. On a typical day, a driver might submit two or three booking requests but check their booking status ten or twenty times. This means the system should be optimised for read performance, and read-path caching is highly valuable.\\

\textbf{Globally distributed but locally clustered users:}\\ Drivers are distributed across the entire world, but within any given region, they are clustered around cities and major road networks. This has two implications. First, the system must be deployed in multiple geographic regions to keep latency acceptable for all users. Second, within a region, road segment capacity contention is highly localised — drivers in Dublin are contending for Dublin road segments, not for segments in Tokyo. This means that conflict detection can be partitioned by geographic region with minimal cross-region coordination. \\

\textbf{Highly skewed data access:}\\
Some road segments are far more frequently accessed than others. Major urban arterials and motorway junctions in large cities will be the subject of tens of thousands of booking requests during morning peak hours. Rural segments may see only a handful of bookings per day. This skew means that a simple even-distribution partitioning strategy may create hot spots on heavily trafficked segments, and the caching strategy must account for this.\\

\textbf{Bursty write patterns:}\\
Booking submissions are not evenly distributed throughout the day. They are heavily concentrated during the periods just before morning and evening peak commute hours, as drivers submit bookings for their upcoming journeys. This means the system must be able to absorb bursts of write traffic that may be five to ten times the average rate, without queuing delays that would cause booking acknowledgments to exceed the 200-millisecond target.\\

\textbf{Failures are rare but must be handled:}\\
Under normal operating conditions, individual service failures are rare. However, at the scale of a globally deployed system, some node will be failing at almost any given moment. The system must handle these failures transparently without requiring manual intervention or accepting data loss.

\section{Specification}

\begin{figure}[H]
\centering
\includegraphics[width=\textwidth]{figures/ds-architecture.png}
\caption{System architecture all services and communication paths}
\label{fig:architecture}
\end{figure}

\subsection{Core Processing Model}
The system enforces a pre-booking policy through an asynchronous pipeline. When a driver submits a booking, the system immediately acknowledges receipt and begins processing the booking in the background. This design choice is deliberate: the full capacity check requires database operations that, under load, could take hundreds of milliseconds. Making the driver wait synchronously for this would make the booking submission endpoint slow and would couple its latency directly to the performance of the database layer. \\
Instead, the booking is given a PENDING status upon receipt. The system then asynchronously checks whether the route's road segments have sufficient capacity for the driver's requested time window. When the check completes, the booking is updated to APPROVED or REJECTED and the driver is notified. From the driver's perspective, they submit a booking, receive an immediate acknowledgement, and then receive a notification within a few seconds. \\
This model is a practical implementation of the Saga pattern: rather than a single distributed transaction spanning multiple services, the booking workflow is decomposed into a sequence of local transactions, each of which either succeeds and triggers the next step or fails and triggers a compensating action.

\subsection{Services}
The system is decomposed into seven services, each with a single well-defined responsibility. All client traffic enters through the API Gateway, which validates JWTs on every request except the public authentication paths, injects an X-Driver-ID header into forwarded requests, and routes based on path prefix. POST /auth/register and POST /auth/login are the only endpoints accessible without a valid JWT. All other endpoints require a Bearer token in the Authorization header.\\\\

\textbf{BFF (Backend for Frontend):}\\
Serves the compiled React frontend to browsers and proxies API calls from the frontend to the API Gateway. Owns no database tables. Exposes no API of its own.\\\\

\textbf{API Gateway:}\\
The single entry point for all external traffic. Validates authentication tokens, enforces rate limits using Redis counters, and routes requests to the appropriate downstream service based on path prefix. Owns no database tables.\\\\

\textbf{Driver Service:}\\
Manages driver identity. Handles registration, authentication, and the issuance of signed session tokens. Owns the drivers table.\\
POST /auth/register accepts a driver's name, email, password, licence number, vehicle type, and home region. It creates the account and returns a JWT. Passwords are stored as bcrypt hashes and the plain-text password is never stored or logged.\\
POST /auth/login accepts an email and password. On success it returns a signed JWT containing the driver's ID and a 24-hour expiry timestamp. This token must be presented on all subsequent authenticated requests.\\
GET /drivers/me returns the authenticated driver's profile information.\\\\

\textbf{Booking Service:}\\
Owns the booking lifecycle. Accepts booking submissions, maintains the booking state machine, and coordinates the event pipeline that drives conflict detection and notification delivery. Owns the bookings table.\\
POST /bookings accepts a route ID and departure time. Returns HTTP 202 Accepted immediately with the new booking's ID. The booking is recorded with PENDING status and enters the asynchronous processing pipeline.\\
GET /bookings returns a paginated list of the authenticated driver's bookings including their current status and route details.\\
GET /bookings/\{booking_id\} returns the full detail of a single booking.\\
DELETE /bookings/\{booking_id\} cancels a booking. Valid for bookings in PENDING or APPROVED status. Returns an error if the booking is in any other state.\\\\

\textbf{Route Service:}\\
Computes and caches driving routes by integrating with the OSRM routing engine. Owns the routes table and the road_segments table for seeding purposes.\\
GET /routes accepts origin and destination coordinates as query parameters. If a route between those coordinates has been computed previously the cached result is returned. Otherwise the service queries OSRM, stores the result, and returns it. The response includes the route geometry as a GeoJSON LineString and the ordered list of road segments with their current utilisation levels.\\
GET /routes/\{route_id\} returns the full detail of a specific route.\\
GET /routes/\{route_id\}/segments returns the ordered segment list for a route with capacity information for each segment.\\\\

\textbf{Conflict Detection Service:}\\
The capacity enforcement engine. Consumes booking events, checks road segment capacity, creates or denies reservations, and emits assessment outcomes. Owns the road_segments and segment_reservations tables for operational purposes. Exposes no external API; it operates entirely through event consumption and publication.\\\\

\textbf{Messaging Service:}\\
Delivers notifications to drivers. Consumes booking outcome events and creates inbox messages. Owns the messages table.\\
GET /messages returns the authenticated driver's notification inbox ordered by most recent first.\\
PUT /messages/\{message_id\}/read marks a specific message as read.

\subsection{Frontend pages}
The React frontend served by the BFF provides the following pages:\\\\
/login presents the driver login form.\\\\
/register presents the driver registration form including all required fields.\\\\
/routes displays an interactive Mapbox map. Drivers can select an origin and destination, view the computed route overlaid on the map, and see road segments colour-coded by capacity utilisation (green for available, red for full). Drivers can select a departure time and submit a booking directly from this page.\\\\
/bookings shows the driver's booking list with current statuses. Drivers can view booking details and cancel bookings from this page. \\\\
/inbox shows the driver's message inbox with read and unread indicators.

\subsection{Internal Event Flows}
Services do not call each other's HTTP endpoints directly. All inter-service coordination uses asynchronous events published to Redis Streams. The three event flows are as follows.\\\\
booking.created is published by the Booking Service when a new booking is persisted. It carries the booking ID, driver ID, route ID, and departure time. It is consumed by the Conflict Detection Service, which uses it to trigger the capacity check. \\\\
route.assessed is published by the Conflict Detection Service after evaluating capacity. It carries the booking ID, route ID, and a boolean indicating whether all segments were available. It is consumed by the Booking Service, which updates the booking status accordingly. \\\\
booking.updated is published by the Booking Service when a booking's status changes. It carries the booking ID, driver ID, and the new status. It is consumed by the Messaging Service, which creates an inbox notification for the driver.

\subsection{Fault Behaviour}
The system is designed to behave predictably in the presence of the failures described in the failure model.\\

\textbf{Redis becomes unavailable:}\\ The API Gateway's rate limiter fails open, meaning all requests are allowed through without rate limiting. This preserves system availability at the cost of rate enforcement during the outage. The Conflict Detection Service bypasses its segment capacity cache and queries PostgreSQL directly. The Booking Service's Transactional Outbox continues accumulating events in PostgreSQL until Redis recovers, at which point the relay resumes publishing and no booking events are lost.\\

\textbf{A service instance crashes:}\\ Kubernetes detects the failed pod via its liveness probe and automatically restarts it. Any event the crashed instance was consuming remains unacknowledged in the Redis Streams consumer group. When the instance restarts it re-reads the unacknowledged event and processes it again. Consumer-side idempotency checks in the Messaging Service ensure that duplicate event deliveries do not produce duplicate notifications.\\

\textbf{The database primary becomes unavailable:}\\ Write operations fail and services return error responses to clients. The system follows a single-primary model and CloudNativePG manages automatic failover in the Kubernetes deployment. During a primary outage, events continue to queue in Redis Streams and are processed once a new primary is elected.\\

\textbf{Two drivers simultaneously request the last available segment slot:}\\
The Conflict Detection Service uses pessimistic row-level locking to serialise these requests at the database level. One transaction proceeds and consumes the last slot. The second transaction blocks, unblocks when the first commits, re-reads the capacity count, finds zero slots remaining, and produces a REJECTED outcome. \\

\textbf{A pod is manually deleted during a demo:}\\ Kubernetes automatically reschedules the pod. Because all services run a minimum of two replicas, the remaining replica continues serving traffic during the restart. The restarted pod rejoins its Redis Streams consumer group and resumes processing.

\subsection{Gaps and known Limitations}
The following items represent gaps between the intended design and the current implementation.\\
The Redis caching layer for segment capacity data in the Conflict Detection Service is designed but not yet implemented. Every capacity check currently queries PostgreSQL directly. This is correct but will become a bottleneck at high load because it forgoes the read optimisation that the cache would provide.\\
Automated testing is not present in the repository. All verification has been performed manually through the frontend and direct API calls.

\section{Architecture and Design}
\subsection{Architectural Style and Rationale}
The system follows a microservices architecture in which the application is decomposed into a set of small, independently deployable services that each own a clearly defined slice of the system's functionality and data. Services communicate through well-defined interfaces rather than sharing memory or internal state directly.\\
This choice was made specifically because of the global scale requirement. A monolithic application cannot be scaled selectively. If the conflict detection logic becomes a bottleneck during peak hours, a monolith forces scaling everything, including the parts that are not under pressure. With microservices, each component can be scaled independently based on actual demand.\\
The architecture also provides fault isolation. A bug or crash in the Messaging Service cannot crash the Booking Service. A slow Route Service does not affect conflict detection. Services can be deployed, updated, and restarted independently without taking the entire system offline.\\
The trade-off is operational complexity. Running seven services instead of one requires more infrastructure and more careful design of interactions between services. This complexity is justified by the scale and availability requirements of the system. The team mitigated this complexity by enforcing a consistent internal structure across all services.

\subsection{Bounded Contexts and Service Ownership}
Each service that manages its own database table represents a bounded context in the domain-driven design sense. A bounded context is a region of the system within which a particular domain model applies and outside of which it does not. The boundaries were drawn along natural domain lines: the Driver Service owns driver identity and authentication, the Booking Service owns the booking lifecycle, the Conflict Detection Service owns segment capacity and reservations, the Route Service owns route computation and road segment data, and the Messaging Service owns the driver inbox.\\
These boundaries mean that each service can evolve its data model and business rules independently. A change to how the Messaging Service stores messages does not require any change to the Booking Service. A change to the Conflict Detection Service's reservation schema does not affect the Driver Service.\\
\begin{figure}[H]
\centering
\includegraphics[width=\textwidth]{figures/data_model.png}
\caption{Domain data model}
\label{fig:datamodel}
\end{figure}

All services share a single PostgreSQL database but each service has its own SQLAlchemy metadata registry. Because SQLAlchemy cannot resolve foreign key references across separate registries, cross-context references are handled by storing foreign IDs as plain typed columns rather than ORM-level foreign keys. The actual referential integrity constraints are created in Alembic migrations at the database level, so the database still enforces correctness even though the ORM layer does not express the relationship directly.

\subsection{Communication Model}
The system uses two distinct communication models chosen based on the nature of each interaction.\\
Synchronous HTTP REST is used for all client-facing interactions including registration, login, route lookup, booking submission, status queries, and inbox access. These interactions require an immediate response. A driver submitting a booking needs to know within 200 milliseconds that their request was received. Synchronous HTTP is appropriate here because the response can be generated quickly without waiting for conflict detection to complete.\\
Asynchronous event streams are used for all inter-service coordination. The booking approval workflow involves multiple services and takes a variable amount of time. Making this workflow synchronous would create tight coupling. If the Conflict Detection Service is slow or temporarily down, a synchronous design would make the booking submission endpoint slow or unavailable as well. By using asynchronous events, the services are temporally decoupled. Each service processes events at its own pace, and a slow or temporarily unavailable service causes events to queue rather than causing failures to propagate upstream to the driver.\\
Redis Streams was chosen as the message broker. It provides persistent, ordered, consumer-group-based event delivery with at-least-once semantics. Events survive Redis restarts via append-only file persistence. Consumer groups allow multiple instances of a service to share event processing load, with Redis ensuring each event is delivered to exactly one instance in the group. All three event streams in the system use consumer groups for this reason.

\subsection{Observability and Health Monitoring}
In a distributed system with multiple services, knowing the health of each component and being able to trace requests across service boundaries is essential for operations and debugging.

Each service exposes three health endpoints following a standard pattern. The /health endpoint provides a basic liveness check that returns a simple "ok" response, used by orchestrators to determine if a container is running. The /health/live endpoint serves as the Kubernetes liveness probe, indicating whether the service process is alive and able to handle requests. The /health/ready endpoint serves as the Kubernetes readiness probe and checks all dependencies required for the service to function. For example, the Conflict Detection Service's readiness check verifies connectivity to both PostgreSQL and Redis. If either dependency is unavailable, the endpoint returns HTTP 503 indicating the service cannot serve traffic.

Request tracing across services is achieved through the X-Correlation-ID header. The API Gateway generates a unique correlation ID for every incoming request, either using an ID provided by the client or generating a fresh UUID. This correlation ID is propagated to all downstream services, allowing operators to trace a single request as it flows through the booking pipeline. When a driver submits a booking, the correlation ID travels from the API Gateway through the Booking Service, Conflict Detection Service, and Messaging Service, appearing in all log entries. This is particularly valuable when investigating issues in production or during demo failure scenarios.

The current implementation uses basic Python logging. Each service logs at the INFO level with format that includes the correlation ID when available. While this provides sufficient visibility for the current scope, a production deployment would benefit from structured logging with OpenTelemetry for distributed tracing and Prometheus for metrics collection. These additions would enable alerting on error rates and latency percentiles across the service fleet.

\subsection{Reliability Patterns}
The asynchronous event-driven architecture introduces several failure modes that must be handled explicitly. This section describes the patterns used to ensure that no booking is lost and no event is processed incorrectly, even in the presence of crashes and restarts.\\\\
\textbf{The Transactional Outbox:}\\
The most subtle reliability problem in an event-driven system is the dual-write problem. When a new booking is created, the Booking Service must both persist the booking to PostgreSQL and publish a booking.created event to Redis Streams. These are two separate systems. If the service writes to PostgreSQL and then crashes before publishing the event, the booking will sit in PENDING forever because the conflict detection pipeline will never be triggered for it.\\
The naive solution of writing to both systems in sequence does not work because there is a window between the two writes during which a crash produces an inconsistent state. The Transactional Outbox pattern solves this cleanly. Instead of writing to Redis directly, the service writes to an outbox_events table in the same PostgreSQL transaction as the business data. The business record and the outbox entry are committed atomically. A separate background process called the relay polls the outbox table for unsent entries, publishes them to Redis Streams, and marks them as sent. The relay polls every 500 milliseconds and processes events in batches of 50. It acquires outbox rows using SELECT ... FOR UPDATE SKIP LOCKED, which prevents multiple relay instances from publishing the same event. Published events are retained for seven days before a cleanup job removes them.\\
If the relay publishes an event to Redis but crashes before marking it as sent, the event will be published again on the next relay cycle. This produces at-least-once delivery to Redis, which is safe because consumers are idempotent. Both the Booking Service and the Conflict Detection Service use this pattern for their outgoing events. In the Booking Service the relay runs as a background goroutine; in the Conflict Detection Service it runs as a background thread.\\\\
\textbf{Idempotent Consumers:}\\
Because the outbox relay provides at-least-once delivery, consumers must tolerate duplicate events without producing duplicate side effects. Each consumer maintains a processed_events table with a composite primary key of (event_id, consumer_name). Before processing an event, the consumer checks whether that event ID has already been recorded. If it has, the event is silently skipped. The insert uses ON CONFLICT DO NOTHING so that even under race conditions between consumer replicas, no duplicate processing occurs. This pattern is used by the Booking Service when consuming route.assessed events and by the Messaging Service when consuming booking.updated events.\\\\
\textbf{Dead Letter Queues:}\\
Events that fail processing after a configurable number of retries are moved to a dead letter queue implemented as a separate Redis Stream named \{stream_name\}.dlq. This prevents a single malformed or unprocessable event from blocking the processing of all subsequent events, a problem known as head-of-line blocking. Events in the dead letter queue can be inspected manually and either corrected and resubmitted or discarded. All three event-consuming services use this pattern.\\\\
\textbf{Booking Expiry:}\\
A background goroutine in the Booking Service runs every 30 seconds and transitions all PENDING or APPROVED bookings whose departure time has passed to EXPIRED. This ensures that stale bookings do not permanently hold segment capacity and that approved journeys which were never actually taken do not consume capacity indefinitely. The 30-second interval is a deliberate trade-off between responsiveness and database load.

\subsection{Consistency and Concurrency Control}
The hard correctness requirement of the system is that no road segment is ever booked beyond its capacity. This section describes how the system enforces this invariant and how the booking lifecycle is modelled to support it.\\\\
\textbf{Booking State Machine:}\\
A booking object exists in exactly one of five states at any point in time. The valid transitions are as follows.\\
PENDING is the initial state. The booking has been received but not yet assessed.\\\\
PENDING transitions to APPROVED when the Conflict Detection Service determines that all segments along the route have sufficient capacity during the booking's time window.\\\\
PENDING transitions to REJECTED when the Conflict Detection Service determines that at least one segment is at capacity.\\\\
PENDING transitions to EXPIRED when the booking's departure time passes without the booking having been resolved. This is enforced by the expiry loop described above.\\\\
APPROVED transitions to CANCELLED when the driver explicitly cancels the booking after it has been approved. This releases the segment reservations that were created during the approval process.\\\\
PENDING can also transition to CANCELLED if the driver cancels before the booking has been assessed.\\\\
Once a booking reaches REJECTED, EXPIRED, or CANCELLED it cannot change state again. These are terminal states.\\\\
\textbf{Pessimistic Locking for Capacity Enforcement:}\\
To prevent two concurrent bookings from both reading stale capacity counts and both being approved when only one slot remains, the Conflict Detection Service uses pessimistic row-level locking via SELECT ... FOR UPDATE in SQLAlchemy using .with_for_update(). This tells PostgreSQL to lock the selected segment rows for the duration of the current transaction. Any other transaction attempting to lock the same rows will block until the first transaction commits or rolls back. The lock granularity is at the individual segment row level, meaning that bookings on non-overlapping routes proceed in parallel without any blocking.\\
This is a deliberate choice of pessimistic over optimistic concurrency control. Optimistic control would be appropriate if conflicts were rare. However, during peak hours many drivers compete for capacity on the same popular segments, making conflicts frequent. Pessimistic locking is more efficient in high-contention scenarios because it avoids the cost of repeated transaction retries.\\\\
\textbf{The Time Window Model:}\\
When the Conflict Detection Service receives a booking.created event it retrieves the list of road segments for the route. It calculates the time window during which the driver will occupy each segment using a staggered model. The driver is estimated to arrive at segment N at departure_time + (N multiplied by 300 seconds). Each segment is reserved for a 300-second window starting at that calculated arrival time. This is more realistic than a flat simultaneous-occupancy assumption and reduces unnecessary capacity contention on segments the driver will not reach until later in the journey.

\subsection{Global Deployment Architecture}
The system is designed for deployment in multiple geographic regions, each containing at least two availability zones. A region is a broad geographic area such as Western Europe, North America East, or Asia Pacific. Each region runs a complete independent instance of all seven services and its own database cluster. Drivers are routed to their nearest region by a global DNS load balancer using latency-based routing. \\
Within a region, road segment data is partitioned by sub-region. A booking for a journey in Dublin is processed entirely within the Western Europe region and requires no communication with other regions. This means that conflict detection, the most latency-sensitive internal operation, is always a local operation within a region.

\subsection{Data Partitioning Strategy}
Road segment capacity data has an important property: it is inherently local. A booking for a journey in Dublin only needs to check capacity on Dublin road segments. A driver in Tokyo never contends for capacity on a segment in Oslo. This means that if the system were deployed at global scale, the segment data could be partitioned by geographic region with minimal cross-partition coordination. Each region's bookings would be resolved entirely within that region's data partition.\\
The data model has been designed with this future partitioning in mind. Every road segment record and every driver record carries a region field that could serve as a partition key. The current deployment does not partition the data. All segments and reservations live in a single CloudNativePG-managed PostgreSQL cluster. At the scale of the current system this is sufficient, and introducing partitioning prematurely would add operational complexity without measurable benefit. The region field is in place so that partitioning can be introduced when load requires it, without schema changes.

\subsection{Caching Strategy}
The system uses caching at multiple layers to reduce latency and database load.\\
Route caching is the most straightforward cache in the system. Computed routes are stored permanently in PostgreSQL the first time they are computed by the Route Service. Routes do not change frequently because the road network is stable. When a driver requests a route between two coordinates that have been queried before, the cached result is returned directly without querying OSRM again. This is a write-once, never-evict cache populated lazily on first access.\\\\
Segment capacity caching is designed to use Redis with a 30-second TTL, with the cache key structured as segment:capacity:\{segment_id\}:\{time_window\}. However, as noted in the gaps section, this cache has not yet been implemented. All capacity checks currently go directly to the database. On Redis failure the design specifies that the service falls through to the database directly, maintaining correctness at the cost of additional database load. \\\\
Rate limit counters are stored in Redis with a 60-second TTL per driver. The API Gateway increments this counter on each request and rejects requests that exceed 100 per minute with an HTTP 429 response. On Redis failure the rate limiter fails open and allows all requests through.\\\\
JWT validation requires no network call or database query. All services receive the shared JWT secret at startup and validate tokens locally. This makes authentication essentially free in terms of latency.

\subsection{Replication Strategy}
Database replication uses a single-primary, multiple-replica model per region. All writes go to the primary. The single-primary model was chosen deliberately over multi-master replication. Multi-master replication would allow writes to be distributed across replicas but introduces the possibility of write conflicts. For a system where the hard correctness requirement is that segment capacity is never overbooked, introducing write conflicts at the database level would add significant complexity and risk. CloudNativePG manages automatic failover in the Kubernetes deployment. \\
Service replication runs all application services as multiple pod replicas behind a Kubernetes Service load balancer. The minimum configuration is two replicas per service. This means a single pod failure does not cause any downtime because the remaining replica continues serving traffic while Kubernetes restarts the failed pod. During peak hours, horizontal pod autoscaling can increase the replica count based on CPU utilisation.\\
Redis replication uses Redis with persistence enabled via RDB snapshots and append-only file logging to ensure events survive restarts without loss.

\subsection{Load Balancing}
At the global level a DNS-based load balancer routes each driver to their nearest regional deployment based on measured network latency. This is transparent to the client and represents the first level of geographic traffic distribution.\\
Within the Kubernetes cluster each service is fronted by a Kubernetes ClusterIP Service that distributes requests across pod replicas using round-robin. This is Layer 4 load balancing provided by kube-proxy.\\
For event processing, Redis Streams consumer groups provide automatic load balancing. When multiple replicas of the Conflict Detection Service are running, Redis delivers each event to exactly one replica, distributing the processing load evenly. This means that scaling the Conflict Detection Service up during peak hours immediately results in proportionally higher event processing throughput with no configuration changes required.

\subsection{Service-Specific Design}
\textbf{Booking Service:}\\
The Booking Service is written in Go. This choice was deliberate: Go's lightweight goroutines allow the service to manage multiple concurrent background tasks alongside the HTTP server within a single process with minimal memory overhead. The three background tasks are the outbox relay, the expiry loop, and the route.assessed event consumer, all running as separate goroutines managed by the main process.\\\\
\textbf{Conflict Detection Service:}\\
The Conflict Detection Service is the most algorithmically critical component in the system. Its primary responsibility is consuming booking.created events, checking segment capacity using the pessimistic locking and time window model described above, creating or denying reservations, and publishing the assessment result via its own transactional outbox. When a booking is cancelled, it also consumes the booking.updated event and releases the corresponding segment reservations.\\\\
\textbf{Route Service and OSRM Integration:}\\
The Route Service integrates with OSRM, a standalone routing engine that runs as a separate container pre-loaded with OpenStreetMap road data for Ireland. OSRM computes the fastest driving path between any two coordinates, returning the route geometry as GeoJSON and an ordered list of OSM way IDs.\\
The Route Service maps these OSM way IDs to the system's internal road segment IDs via the osm_way_id column on the road_segments table. This mapping translates a human-meaningful route suitable for rendering on a Mapbox map into a machine-meaningful sequence of segment IDs that the Conflict Detection Service can use to check capacity.\\
No other service communicates with OSRM directly. The Route Service is the sole owner of the OSRM integration.\\
Road segment data is seeded into the database using the Overpass API, a query interface for OpenStreetMap data. The seeding script queries Overpass for roads of specified types within a geographic bounding box, extracts way IDs, names, and coordinates, and inserts them into the road_segments table.

\section{Implementation}
\subsection{Technology Stack}
The system uses a polyglot technology stack with each technology chosen for specific reasons. \\
Python with FastAPI is used for the API Gateway, Driver Service, Route Service, Conflict Detection Service, and Messaging Service. FastAPI was chosen for its developer ergonomics, automatic API documentation generation, and performance relative to other Python web frameworks. It uses Python's asyncio for non-blocking I/O, which makes it well suited to services that spend significant time waiting on database or Redis responses.\\
Go with the chi router is used for the Booking Service. Go's native concurrency model makes it straightforward to run the HTTP server, the outbox relay goroutine, and the expiry loop goroutine concurrently within a single lightweight process.\\
React 19 with TypeScript is used for the frontend. The frontend uses Zustand for authentication state management persisted to localStorage, and TanStack React Query for server state management with automatic background refetching. Mapbox GL JS renders the interactive map, route overlays, and segment capacity colour coding.
PostgreSQL is the primary data store for all persistent data. \\
PostgreSQL was chosen because the booking workflow requires ACID transactional guarantees. The segment capacity check using SELECT ... FOR UPDATE is a SQL-native operation. Alembic manages all schema migrations.\\

Redis handles three distinct responsibilities: the message broker via Redis Streams, the rate limit counter store, and the segment capacity cache (designed but not yet implemented). \\

Docker and Kubernetes provide the container runtime and orchestration layer. The local development environment uses Docker Compose. The demo environment uses k3s running across multiple team laptops to form a multi-node cluster that can demonstrate node failure scenarios. \\
OSRM runs as a standalone container pre-loaded with OpenStreetMap data for Ireland. It is an internal dependency of the Route Service and is not accessible to other services or to external clients.

\subsection{Consistent Service Structure}
main.py contains the FastAPI application instance and router registration. It is intentionally thin, containing no business logic.\\\\
database.py contains the SQLAlchemy engine setup and the session dependency used by router functions.
models.py contains the SQLAlchemy table model definitions for the tables owned by that service.\\\\
schemas.py contains the Pydantic request and response models that define the shapes of data entering and leaving the service's HTTP endpoints.\\\\
router.py contains the HTTP endpoint handler functions. These are intentionally thin, delegating all business logic to the service layer.\\\\
service.py contains the business logic.\\\\
consumer.py contains the Redis Streams event consumer loop for services that consume events.

\subsection{Cross Service Foreign Keys}
All services share a single PostgreSQL database instance but each service defines its own SQLAlchemy declarative_base(). This means SQLAlchemy's ORM layer has no knowledge of tables defined in another service's models. Attempting to use a standard SQLAlchemy ForeignKey('other_table.id') reference across service boundaries causes a NoReferencedTableError at startup because the referenced table is not in the same metadata registry.\\
The solution used throughout the system is to declare cross-service reference columns as plain typed columns in SQLAlchemy, for example mapped_column(UUID(as_uuid=True), nullable=False), and to define the actual database-level foreign key constraint in the Alembic migration file rather than in the model. This means the database enforces referential integrity correctly while the SQLAlchemy ORM layer for each service only knows about its own tables.

\subsection{End to End Booking flow}
The complete lifecycle of a booking from submission to notification involves four services and three Redis Stream events. \\
The driver submits a POST request to /bookings through the API Gateway. The gateway validates the JWT, checks the driver's rate limit counter in Redis, and forwards the request to the Booking Service with an injected X-Driver-ID header. \\
The Booking Service receives the request and inserts a new booking record with PENDING status and a corresponding outbox_events entry into PostgreSQL within a single database transaction. It returns HTTP 202 Accepted to the client. \\
The outbox relay goroutine wakes on its polling interval, reads the pending outbox entry, publishes the booking.created event to Redis Streams, and marks the entry as published. \\
The Conflict Detection Service's consumer loop receives the event via XREADGROUP. It fetches the route's segment list, calculates staggered time windows for each segment, acquires SELECT ... FOR UPDATE locks on the relevant segment rows, counts existing reservations against each segment's capacity, creates reservations if all segments have room, and emits a route.assessed event. It then acknowledges the booking.created event to the consumer group. \\
The Booking Service's consumer loop receives the route.assessed event, verifies the booking is still in PENDING status, updates the status to APPROVED or REJECTED, and publishes a booking.updated event. \\

The Messaging Service's consumer loop receives the booking.updated event, checks its processed_events table to confirm the event has not already been handled, creates an inbox message for the driver, and records the event ID as processed. \\
The driver's frontend, which polls booking status every 5 seconds and the inbox every 10 seconds, displays the updated status and the notification.

\subsection{Concurrent Booking Resolution}
When two drivers simultaneously request the last available slot on a road segment the following sequence occurs.\\
Both drivers submit bookings. Both are recorded as PENDING. Both booking.created events are published to Redis Streams. If two instances of the Conflict Detection Service are running, Redis delivers one event to each instance. Both instances begin processing simultaneously and both issue SELECT ... FOR UPDATE queries on the same segment row.\\
PostgreSQL grants the lock to whichever transaction acquires it first. The second transaction is blocked at the lock acquisition step and waits. \\
The first transaction reads the reservation count, finds one slot remaining, creates the reservation, emits a route.assessed event with APPROVED status, and commits. The lock is released. \\
The second transaction unblocks, reads the reservation count, finds zero slots remaining, does not create a reservation, emits a route.assessed event with REJECTED status, and commits. \\
Driver A's booking is updated to APPROVED and Driver B's booking is updated to REJECTED. Both drivers receive inbox notifications. Driver B is advised to try a different departure time or route.

\subsection{Idempotent Event Processing}
Because the system uses at-least-once event delivery, consumers must handle duplicate deliveries without producing duplicate side effects.\\
The Messaging Service achieves this through a processed_events table. Before processing any event, the service attempts to insert the event's unique ID into this table. If the insert succeeds the event is new and should be processed. If it fails with a unique constraint violation the event is a duplicate and is silently skipped. This converts at-least-once delivery into effectively exactly-once processing for notification creation.\\
The Booking Service handles duplicate route.assessed events through its state machine. The status update query only matches bookings in PENDING status. A booking that has already been updated to APPROVED or REJECTED will not be matched, so duplicate processing has no effect.


\subsection{Failure Handling}
\textbf{Redis failure:}\\
When Redis becomes unavailable the rate limiter in the API Gateway fails open, allowing all requests to proceed. The Conflict Detection Service bypasses the segment capacity cache and falls through to direct database queries. The Booking Service's outbox relay pauses because it cannot publish to Redis Streams, but events continue to accumulate safely in the PostgreSQL outbox table. When Redis recovers the relay resumes publishing the accumulated entries. This behaviour is explicitly designed into the system as documented in the project specification.\\

\textbf{Service instance crash:}\\
When a service pod crashes, Kubernetes detects the failure via the /health liveness probe on each service and automatically restarts the pod. Events that were being processed at the time of the crash remain unacknowledged in the Redis Streams consumer group. When the pod restarts it rejoins the consumer group and re-reads the unacknowledged events. The Messaging Service's processed_events table ensures that a redelivered booking.updated event does not produce a duplicate notification. The Booking Service's state machine ensures that a redelivered route.assessed event has no effect on a booking that has already left PENDING status.
\\\\
\textbf{Database partition:}\\
The system follows a single-primary model. Write operations are only accepted by the partition containing the PostgreSQL primary, preventing split-brain inconsistencies. Services on the side of the partition without the primary return errors for write operations but can continue serving reads. When the partition heals, the system resumes normal operation.\\\\

\textbf{Concurrent conflicting bookings:}\\
Two drivers booking the last available slot simultaneously is handled entirely by the pessimistic locking mechanism described in Section 4.5. This scenario has been tested manually and verified to produce exactly one approval and one rejection consistently.

\section{Deployment}
The system is designed to run in containerised environments using Kubernetes for orchestration. This section describes both the local development setup and the production EKS deployment.

\subsection{Development Environment}
During development the entire system runs locally using Docker Compose. Each service is defined as a container with its own port exposed to the host machine. The docker-compose.yml file defines the service dependencies, ensuring that the API Gateway starts only after all backend services are ready. PostgreSQL and Redis are started as supporting containers, and OSRM runs as a separate container with pre-loaded OpenStreetMap data for Ireland. The frontend is built with Vite and the compiled assets are served by the BFF service. This setup allows developers to make changes to any service and quickly test the full booking flow without requiring cloud infrastructure.

The Helm charts in the /charts directory provide the Kubernetes manifests for production deployment. The charts follow a template-based approach with values files that allow configuration of replica counts, image tags, and resource limits for each service. All service deployments define container resource requests and limits to ensure predictable scheduling behaviour on the Kubernetes cluster.

\subsection{Amazon EKS Production Deployment}
The system is deployed on Amazon EKS with six worker nodes. Each service runs as a Kubernetes Deployment with configurable replica counts. The Kubernetes cluster provides fault tolerance at the infrastructure level: if a worker node fails, the kubelet on that node stops reporting heartbeats and Kubernetes automatically reschedules the pods onto healthy nodes.

PostgreSQL is managed by CloudNativePG, a Kubernetes operator that provides automatic failover and high availability. The operator maintains a primary-replica configuration and automatically promotes a replica to primary if the primary becomes unavailable. This means that a database primary failure does not cause the booking system to become unavailable; the operator promotes a new primary within seconds and services reconnect automatically.

Redis runs as a Deployment with persistence enabled. The append-only file configuration ensures that events published to Redis Streams survive restarts. In a production setting, Redis could be replaced with Amazon ElastiCache for managed Redis with multi-AZ replication for higher availability.

The API Gateway is exposed as a Kubernetes Service of type LoadBalancer, distributing incoming traffic across all API Gateway pods. Each downstream service is similarly exposed via ClusterIP services that load-balance across their pod replicas. Kubernetes liveness and readiness probes use the service health endpoints to ensure traffic is only routed to healthy instances.

Environment configuration is managed through Kubernetes ConfigMaps. The traffic-config ConfigMap contains environment variables such as DATABASE_URL, REDIS_URL, and service URLs for all downstream services. This allows operators to modify configuration without rebuilding container images.

The system follows a single-region deployment model in EKS. While the architecture supports multi-region deployment for global latency optimisation, the current deployment runs in a single AWS region. Drivers connect to the nearest EKS cluster through the API Gateway's public endpoint.

\subsection{Key Design Decisions}
The choice of Kubernetes as the orchestration layer was driven by its industry adoption and the need for automated scaling and failure recovery. The team considered serverless alternatives but decided against them because the background consumer processes require long-running containers rather than short-lived function invocations.

CloudNativePG was chosen for database high availability because it integrates natively with Kubernetes. The operator handles primary election, failover, and backup management without requiring external tools. This keeps the operational footprint smaller than managing a separate PostgreSQL cluster outside of Kubernetes.

Helm was used for templating Kubernetes manifests to avoid duplication across services. Each service deployment follows the same pattern with only the image name, port, and environment variables differing. This makes it easy to add new services or modify existing ones.

\section{Testing}
\subsection{Test Strategy}
The testing strategy addresses three distinct concerns: functional correctness, concurrency integrity, and resilience under failure.

\subsection{Functional testing}
The complete happy-path flow was tested: driver registration, login, route lookup for a Dublin coordinate pair, booking submission, monitoring the booking through the PENDING to APPROVED transition, and verifying inbox notification delivery.\\
Cached route lookups return in under 50 milliseconds. Cold route lookups requiring OSRM computation are the primary performance bottleneck, with latency that varies depending on OSRM backend response time. The transition from PENDING to APPROVED completes within a few seconds under normal load, which is within the notification delivery target.\\
Booking cancellation was tested for PENDING bookings. The APPROVED to CANCELLED transition is specified in the state machine and supported by the implementation.\\
The admin dashboard page displays live service health and booking throughput metrics.

\subsection{Concurrency Testing}
The double-booking prevention test was performed by pre-loading segment reservations so that a target segment had exactly one remaining slot, then simultaneously submitting two booking requests for routes traversing that segment with the same departure time.\\
Atomic transactions on segment reservations hold under concurrent load, with exactly one booking approved and one rejected. This confirms that the pessimistic locking mechanism works correctly.
\\
\subsection{Resilience Testing}
The following scenarios were tested as described in the project specification:\\
Service replica failure was tested by deleting pods using kubectl delete pod. Kubernetes automatically restarted the pods and the system resumed normal operation.\\\\
Redis failure was tested by taking the Redis container down. The Conflict Detection Service fell back to direct database queries and events queued in the Transactional Outbox. When Redis was restored, event processing resumed correctly. \\\\
Data persistence was verified by performing a full system restart. Bookings and messages survived the restart, confirming that PostgreSQL persistent volumes are working correctly.\\\\
Concurrent booking conflict was verified by submitting two simultaneous booking requests for the last available slot on a segment. One booking was approved and one was rejected with no double booking occurring.

\subsection{Test Summary}
\begin{table}[h!]
\centering
\begin{tabular}{|p{6cm}|c|p{7cm}|}
\hline
\textbf{Test Scenario} & \textbf{Result} & \textbf{Notes} \\
\hline
Driver registration and login & Pass & JWT issued and validated correctly \\
\hline
Route lookup cached & Pass & Returns under 50ms \\
\hline
Route lookup uncached & Pass & Latency varies with OSRM backend \\
\hline
Booking submission and approval & Pass & PENDING to APPROVED within a few seconds \\
\hline
Booking cancellation from PENDING & Pass & Correctly transitions to CANCELLED \\
\hline
Concurrent double-booking prevention & Pass & One approval, one rejection, no double booking \\
\hline
Booking rejection notification & Pass & Rejected driver notified correctly \\
\hline
Data persistence across restart & Pass & All data survives full system restart \\
\hline
Redis failure and recovery & Pass & Outbox holds events, processed on recovery \\
\hline
Service replica failure via kubectl & Pass & Kubernetes restarts pod automatically \\
\hline
Rate limit fail-open & Pass & Requests allowed through when Redis unavailable \\
\hline
\end{tabular}
\caption{System Test Results}
\label{tab:test_results}
\end{table}

\section{Allocation of work}
The project was developed by a team of six members. Every member took responsibility of at least one service with all the members contributing to system integration, testing and this report.

\begin{table}[H]
\centering
\begin{tabular}{|p{3.5cm}|p{4cm}|p{7.5cm}|}
\hline
\textbf{Team Member} & \textbf{Primary Service} & \textbf{Key Contributions} \\
\hline
Navina & Booking Service (Go) & State machine design, Transactional Outbox, relay goroutine, expiry loop, cancellation logic \\
\hline
Aditya & Conflict Detection Service & Pessimistic locking, staggered time window model, Dead Letter Queue \\
\hline
Lukas & API Gateway & JWT validation middleware, Redis rate limiting, path-based routing, fail-open degradation \\
\hline
Lukas & Route Service & OSRM integration, route resolution flow, Overpass API seeding script, segment data model, route caching \\
\hline
Shravani & Driver Service & Account registration, bcrypt password hashing, JWT issuance, driver profile management \\
\hline
Pranjal & Messaging Service & Notification creation, idempotent event processing via processed\_events table, inbox management \\
\hline
Lukas, Michal & BFF and Infrastructure & Mapbox integration, Docker Compose, Kubernetes manifests, EKS cluster setup \\
\hline
Lukas, Aditya & Frontend & React frontend including all pages and dashboard \\
\hline
\end{tabular}
\caption{Team Member Contributions}
\label{tab:team_contributions}
\end{table}

\section{Conclusion}
The system implements a working end-to-end booking pipeline across seven microservices. A driver can register, authenticate, look up a route on a real map, submit a booking, and receive a notification with the outcome. The asynchronous event-driven architecture, the Transactional Outbox for reliable event publishing, pessimistic locking for correct capacity enforcement, and dead letter queues for fault tolerance are all implemented and functional in the Docker Compose deployment.\\
The main gaps are in the production infrastructure layer. The Redis caching tier for segment capacity is designed but not yet implemented, and automated testing is absent. These are the areas that would need to be addressed before the system could be considered production-ready, but the core distributed systems concerns of consistency, reliability, and fault tolerance are addressed in the current implementation.

\section{Signed by members:}
Aditya Kapse\\
Lukas Linss\\
Michal\\
Navina\\
Pranjal\\
Shravani Chougule

\end{document}

